% SGA 3, Expose VI_B, Loop 1 English reconstruction.
% Scope: Proposition 6.5.1 in full and source editor's note 61.
% Authority: Polo--Gille Exp6B-13oct24.pdf, component-local page
% 45 (combined-reader page 386; re-edition header 375;
% original-margin printed page 361).
% Theorem 6.5.2 begins on local page 46 / combined page 387
% and is excluded.
% OCR and Reinhold English Markdown are comparison witnesses only.
% This unit contains no diagrams.

\medskip
\phantomsection\label{prop:sga3-VIB-6.5.1}%
\noindent\textbf{命题 6.5.1.---}
\emph{设 \(G\) 为在 \(S\) 上局部有限呈示
且拟分离的 \(S\)-群, \(H\) 为纤维连通的平滑 \(S\)-群,
\(u:H\to G\) 为单态射.
设 \(N\) 为 \(H\) 在 \(G\) 中的正规化子
(参见 \hyperref[def:sga3-VIB-6.1]{6.1}). 由下面的
\hyperref[cor:sga3-VIB-6.5.5]{6.5.5}, \(N\) 可由
\(G\) 的闭子群概形表出, 且在 \(G\) 上有限呈示.
在此情形下, 下列条件
等价:}
\begin{enumerate}
\item[\textup{\SGAInlineRefTarget{prop:sga3-VIB-6.5.1:item:i}(i)}]\phantomsection\label{cond:sga3-VIB-6.5.1-i}
  \emph{典范态射 \(H\to N\) 是开浸入.}
\item[\textup{\SGAInlineRefTarget{prop:sga3-VIB-6.5.1:item:ii}(ii)}]\phantomsection\label{cond:sga3-VIB-6.5.1-ii}
  \emph{\(N^0=H\)
  (参见 \hyperref[thm:sga3-VIB-3.10]{3.10}).}
\item[\textup{\SGAInlineRefTarget{prop:sga3-VIB-6.5.1:item:iii}(iii)}]\phantomsection\label{cond:sga3-VIB-6.5.1-iii}
  \emph{对每个 \(s\in S\), 有
  \(H_s=(N_s)^0\).}
\end{enumerate}

由引理~\hyperref[lem:sga3-VIB-3.10.1]{3.10.1},
条件 \textup{\SGARefLink{prop:sga3-VIB-6.5.1:item:i}{(i)}} 蕴涵 \textup{\SGARefLink{prop:sga3-VIB-6.5.1:item:ii}{(ii)}}, 因为 \(H^0=H\).
显然 \textup{\SGARefLink{prop:sga3-VIB-6.5.1:item:ii}{(ii)}} 也蕴涵 \textup{\SGARefLink{prop:sga3-VIB-6.5.1:item:iii}{(iii)}}, 因为
\[
  H_s=(N^0)_s=(N_s)^0.
\]

最后证明 \textup{\SGARefLink{prop:sga3-VIB-6.5.1:item:iii}{(iii)}} 蕴涵 \textup{\SGARefLink{prop:sga3-VIB-6.5.1:item:i}{(i)}}.
由于对每个 \(s\in S\) 有 \(H_s=(N^0)_s\), 态射
\[
  H_s\longrightarrow N_s
\]
为开浸入. 此外, \(H\) 和 \(N\) 在 \(S\) 上局部
有限呈示, 且 \(H\) 在 \(S\) 上平坦.
因此, \(H\to N\) 为开浸入
(EGA~IV\(_4\), 17.9.5).

为便于读者查阅, 下面收入了
第 XI 篇的结果 \SGARefLink{sga3:xi:theorem:6:8}{6.8}--6.11.\footnote{%
\label{note:sga3-VIB-61}编者注: 我们在此插入了摘自
\SGARefLink{sga3:xi:theorem:6:8}{第 XI 篇, 6.8}--6.11 的 \hyperref[thm:sga3-VIB-6.5.2]{6.5.2}--\SGARefLink{cor:sga3-VIB-6.5.5}{6.5.5} 节.
事实上, A.\ Grothendieck 在 XI~6 开头的一则注中已经
提出这一建议: ``本节不使用
[XI 的] \SGARefLink{sga3:xi:section:3}{第 3 节}、第 4 节和第 5 节的结果;
它自然应当置于
VI\textsubscript{B} 中.''}

% 翻译统计: 源行数 58; 输出行数 61 (含本统计行); 行内/行间数学片段 28; 翻译占位 0. 底本疑点另报主线程.
